\chapter{Generalisierung und Validierung}

\section{Ziel dieses Kapitels}

In den bisherigen Kapiteln haben wir verschiedene Modelle und Trainingsmethoden betrachtet:
lineare Regression (linear regression), logistische Regression (logistic regression), Perzeptrons (perceptrons) und mehrschichtige Perzeptrons (multilayer perceptrons).
Dabei ging es oft darum, eine Fehlerfunktion auf den Trainingsdaten zu minimieren.

Das eigentliche Ziel im maschinellen Lernen ist aber nicht, auf den Trainingsdaten möglichst gut zu sein.
Das eigentliche Ziel ist Generalisierung (generalization).

\begin{conceptbox}
Generalisierung (generalization) bedeutet, dass ein Modell auf neuen, bisher nicht gesehenen Daten gut funktioniert.

Ein Modell soll also nicht nur die Trainingsdaten erklären, sondern die zugrunde liegende Struktur lernen.
\end{conceptbox}

Nach diesem Kapitel solltest du folgende Fragen beantworten können:

\begin{itemize}
    \item Was ist Generalisierung (generalization)?
    \item Was ist der Unterschied zwischen Trainingsfehler (training error), Validierungsfehler (validation error) und Testfehler (test error)?
    \item Was ist Überanpassung (overfitting)?
    \item Was ist Unteranpassung (underfitting)?
    \item Was ist Modellkomplexität (model complexity)?
    \item Was ist der Bias-Varianz-Trade-off (bias-variance trade-off)?
    \item Warum braucht man Validierungsdaten (validation data)?
    \item Was ist Kreuzvalidierung (cross-validation)?
    \item Was sind Hyperparameter (hyperparameters)?
    \item Wie funktioniert Modellselektion (model selection)?
    \item Was ist Regularisierung (regularization)?
    \item Was ist Early Stopping?
    \item Warum darf man Testdaten nicht zum Modellauswählen verwenden?
\end{itemize}

\section{Trainingsfehler und Generalisierung}

Beim Training minimieren wir meistens eine Fehlerfunktion auf den Trainingsdaten:

\[
    E_{\mathrm{train}}(\theta)
    =
    \frac{1}{N_{\mathrm{train}}}
    \sum_{n=1}^{N_{\mathrm{train}}}
    E_n(\theta).
\]

Dabei bezeichnet \(\theta\) allgemein alle Parameter des Modells.
Bei linearer Regression kann \(\theta=\vec{w}\) sein.
Bei neuronalen Netzen enthält \(\theta\) alle Gewichte und Biases.

\begin{definitionbox}
Der Trainingsfehler (training error) ist der Fehler eines Modells auf den Daten, die zum Lernen der Parameter verwendet wurden.
\end{definitionbox}

Ein kleiner Trainingsfehler ist notwendig, aber nicht ausreichend.
Ein Modell kann auf den Trainingsdaten sehr gut sein und trotzdem auf neuen Daten schlecht funktionieren.

\begin{conceptbox}
Das Ziel ist nicht:

\[
    \text{minimiere nur den Trainingsfehler}.
\]

Das Ziel ist:

\[
    \text{minimiere den erwarteten Fehler auf neuen Daten}.
\]
\end{conceptbox}

\section{Daten erzeugende Verteilung}

In vielen theoretischen Betrachtungen nehmen wir an, dass Daten aus einer unbekannten Verteilung stammen.
Ein Datenpunkt besteht aus Eingabe und Zielwert:

\[
    (\vec{x},t).
\]

Diese Paare werden aus einer gemeinsamen Verteilung gezogen:

\[
    (\vec{x},t)
    \sim
    p(\vec{x},t).
\]

Die Trainingsdaten sind dann Stichproben:

\[
    \mathcal{D}_{\mathrm{train}}
    =
    \{(\vec{x}_1,t_1),\dots,(\vec{x}_N,t_N)\}.
\]

Meistens nimmt man an:

\[
    (\vec{x}_1,t_1),\dots,(\vec{x}_N,t_N)
    \overset{\mathrm{i.i.d.}}{\sim}
    p(\vec{x},t).
\]

\begin{notebox}
Die wahre Verteilung \(p(\vec{x},t)\) kennen wir normalerweise nicht.
Wir sehen nur endlich viele Beispiele aus dieser Verteilung.
Genau deshalb ist Generalisierung schwierig.
\end{notebox}

\section{Erwarteter Fehler}

Der ideale Fehler eines Modells wäre der erwartete Fehler über die wahre Datenverteilung.
Für eine Verlustfunktion \(L(y(\vec{x}),t)\) ist der erwartete Fehler:

\[
    E_{\mathrm{true}}(\theta)
    =
    \mathbb{E}_{(\vec{x},t)\sim p(\vec{x},t)}
    \left[
        L(y(\vec{x},\theta),t)
    \right].
\]

Als Integral:

\[
    E_{\mathrm{true}}(\theta)
    =
    \int
    L(y(\vec{x},\theta),t)
    p(\vec{x},t)
    \,d\vec{x}\,dt.
\]

\begin{definitionbox}
Der wahre Generalisierungsfehler (true generalization error) ist der erwartete Fehler eines Modells auf neuen Daten aus der gleichen Datenverteilung:

\[
    E_{\mathrm{true}}(\theta)
    =
    \mathbb{E}
    \left[
        L(y(\vec{x},\theta),t)
    \right].
\]
\end{definitionbox}

Da \(p(\vec{x},t)\) unbekannt ist, können wir \(E_{\mathrm{true}}(\theta)\) nicht exakt berechnen.
Wir müssen ihn mit Daten schätzen.

\section{Trainingsfehler als Schätzer}

Der Trainingsfehler ist ein empirischer Durchschnitt:

\[
    E_{\mathrm{train}}(\theta)
    =
    \frac{1}{N}
    \sum_{n=1}^{N}
    L(y(\vec{x}_n,\theta),t_n).
\]

Wenn \(\theta\) fest wäre und nicht aus denselben Daten gelernt würde, wäre dieser Durchschnitt ein natürlicher Schätzer des erwarteten Fehlers.

Das Problem ist aber:
\(\theta\) wird genau so gewählt, dass \(E_{\mathrm{train}}(\theta)\) klein wird.

\begin{conceptbox}
Der Trainingsfehler ist normalerweise optimistisch verzerrt.

Er unterschätzt häufig den Fehler auf neuen Daten, weil das Modell direkt auf diese Trainingsdaten angepasst wurde.
\end{conceptbox}

Deshalb brauchen wir getrennte Daten zur Bewertung.

\section{Trainings-, Validierungs- und Testdaten}

Typischerweise teilt man die verfügbaren Daten in drei Mengen:

\[
    \mathcal{D}
    =
    \mathcal{D}_{\mathrm{train}}
    \cup
    \mathcal{D}_{\mathrm{val}}
    \cup
    \mathcal{D}_{\mathrm{test}}.
\]

\begin{definitionbox}
Trainingsdaten (training data) werden verwendet, um die Modellparameter zu lernen.

Validierungsdaten (validation data) werden verwendet, um Hyperparameter zu wählen und Modelle zu vergleichen.

Testdaten (test data) werden erst ganz am Ende verwendet, um die endgültige Modellleistung zu schätzen.
\end{definitionbox}

\subsection{Trainingsfehler}

\[
    E_{\mathrm{train}}
    =
    \frac{1}{N_{\mathrm{train}}}
    \sum_{n\in\mathcal{D}_{\mathrm{train}}}
    L(y(\vec{x}_n,\theta),t_n).
\]

\subsection{Validierungsfehler}

\[
    E_{\mathrm{val}}
    =
    \frac{1}{N_{\mathrm{val}}}
    \sum_{n\in\mathcal{D}_{\mathrm{val}}}
    L(y(\vec{x}_n,\theta),t_n).
\]

\subsection{Testfehler}

\[
    E_{\mathrm{test}}
    =
    \frac{1}{N_{\mathrm{test}}}
    \sum_{n\in\mathcal{D}_{\mathrm{test}}}
    L(y(\vec{x}_n,\theta),t_n).
\]

\begin{examtipbox}
Die Testdaten dürfen nicht zur Modellselektion verwendet werden.

Wenn man wiederholt Modelle auf dem Testdatensatz bewertet und danach Entscheidungen trifft, wird der Testdatensatz indirekt zum Validierungsdatensatz.
Dann ist der Testfehler keine ehrliche Schätzung der Generalisierung mehr.
\end{examtipbox}

\section{Überanpassung}

Überanpassung (overfitting) tritt auf, wenn ein Modell die Trainingsdaten sehr gut beschreibt, aber auf neuen Daten schlecht ist.

Typisches Verhalten:

\[
    E_{\mathrm{train}}
    \text{ ist klein},
\]

aber

\[
    E_{\mathrm{test}}
    \text{ ist groß}.
\]

\begin{definitionbox}
Überanpassung (overfitting) bedeutet, dass ein Modell zu stark an die Trainingsdaten angepasst ist und dabei auch Rauschen oder zufällige Besonderheiten der Trainingsdaten lernt.
\end{definitionbox}

\begin{examplebox}
Ein hochgradiges Polynom kann durch viele Trainingspunkte fast exakt verlaufen.
Zwischen den Trainingspunkten kann es aber stark oszillieren.
Dann ist der Trainingsfehler klein, aber der Fehler auf neuen Punkten groß.
\end{examplebox}

\begin{conceptbox}
Überanpassung bedeutet:

\[
    \text{Modell lernt nicht nur Signal, sondern auch Rauschen}.
\]
\end{conceptbox}

\section{Unteranpassung}

Unteranpassung (underfitting) ist das Gegenteil.
Das Modell ist zu einfach, um die Struktur der Daten zu erfassen.

Typisches Verhalten:

\[
    E_{\mathrm{train}}
    \text{ ist groß}
\]

und

\[
    E_{\mathrm{test}}
    \text{ ist ebenfalls groß}.
\]

\begin{definitionbox}
Unteranpassung (underfitting) bedeutet, dass ein Modell zu wenig flexibel ist, um die relevante Struktur in den Daten zu lernen.
\end{definitionbox}

\begin{examplebox}
Wenn die wahre Beziehung gekrümmt ist, aber wir nur eine Gerade verwenden, kann das Modell die Daten nicht gut beschreiben.
Dann liegt Unteranpassung vor.
\end{examplebox}

\section{Modellkomplexität}

Modellkomplexität (model complexity) beschreibt, wie flexibel ein Modell ist.

Beispiele für Komplexität:

\begin{itemize}
    \item Grad eines Polynoms,
    \item Anzahl der Basisfunktionen,
    \item Anzahl der Parameter,
    \item Tiefe eines neuronalen Netzes,
    \item Anzahl der Neuronen,
    \item Stärke der Regularisierung.
\end{itemize}

\begin{conceptbox}
Mit steigender Modellkomplexität sinkt typischerweise der Trainingsfehler.

Der Testfehler sinkt aber nur bis zu einem gewissen Punkt.
Danach kann er wieder steigen, weil das Modell überanpasst.
\end{conceptbox}

Typisches Verhalten:

\[
    \text{zu einfach}
    \quad
    \Rightarrow
    \quad
    \text{Unteranpassung}.
\]

\[
    \text{passend komplex}
    \quad
    \Rightarrow
    \quad
    \text{gute Generalisierung}.
\]

\[
    \text{zu komplex}
    \quad
    \Rightarrow
    \quad
    \text{Überanpassung}.
\]

\section{Beispiel: Polynomiale Regression}

Wir betrachten polynomiale Regression mit Basisfunktionen

\[
    \phi_j(x)=x^j.
\]

Das Modell ist:

\[
    y(x,\vec{w})
    =
    \sum_{j=0}^{M}
    w_jx^j.
\]

Der Polynomgrad \(M\) steuert die Modellkomplexität.

\subsection{Kleines \(M\)}

Wenn \(M=1\), ist das Modell eine Gerade:

\[
    y(x,\vec{w})
    =
    w_0+w_1x.
\]

Das kann für gekrümmte Daten zu einfach sein.

\subsection{Mittleres \(M\)}

Ein mittlerer Polynomgrad kann die Struktur gut erfassen, ohne zu stark zu oszillieren.

\subsection{Großes \(M\)}

Ein sehr großer Polynomgrad kann viele Trainingspunkte exakt treffen, aber zwischen den Punkten stark schwanken.

\begin{conceptbox}
Der Polynomgrad \(M\) ist ein Hyperparameter.

Er wird nicht direkt durch Gradientenabstieg gelernt, sondern muss durch Modellselektion gewählt werden.
\end{conceptbox}

\section{Hyperparameter}

Parameter und Hyperparameter sind unterschiedliche Dinge.

\begin{definitionbox}
Parameter (parameters) werden während des Trainings aus den Trainingsdaten gelernt.

Hyperparameter (hyperparameters) werden vor oder außerhalb des eigentlichen Trainings festgelegt.
\end{definitionbox}

Beispiele für Parameter:

\[
    \vec{w}
\]

bei linearer Regression,

\[
    W^{(l)},\vec{b}^{(l)}
\]

bei neuronalen Netzen.

Beispiele für Hyperparameter:

\begin{itemize}
    \item Polynomgrad \(M\),
    \item Regularisierungsstärke \(\lambda\),
    \item Lernrate \(\eta\),
    \item Anzahl der versteckten Schichten,
    \item Anzahl der Neuronen pro Schicht,
    \item Batch-Größe,
    \item Anzahl der Trainings-Epochen,
    \item Wahl der Aktivierungsfunktion.
\end{itemize}

\begin{notebox}
Hyperparameter werden typischerweise mit Validierungsdaten gewählt.
Wenn man sie mit Testdaten wählt, ist der Testfehler nicht mehr unabhängig.
\end{notebox}

\section{Validierung}

Validierung (validation) bedeutet, Modelle oder Hyperparameter auf Daten zu bewerten, die nicht für das Fitten der Parameter verwendet wurden.

Vorgehen:

\begin{enumerate}
    \item Teile die Daten in Trainings-, Validierungs- und Testdaten.
    \item Trainiere verschiedene Modelle auf den Trainingsdaten.
    \item Berechne den Validierungsfehler.
    \item Wähle das Modell mit dem besten Validierungsfehler.
    \item Bewerte das gewählte Modell einmal auf den Testdaten.
\end{enumerate}

\begin{definitionbox}
Modellselektion (model selection) ist die Auswahl eines Modells oder von Hyperparametern anhand von Validierungsleistung.
\end{definitionbox}

\begin{conceptbox}
Trainingsdaten beantworten:

\[
    \text{Welche Parameter soll das Modell haben?}
\]

Validierungsdaten beantworten:

\[
    \text{Welches Modell oder welche Hyperparameter sollen gewählt werden?}
\]

Testdaten beantworten:

\[
    \text{Wie gut ist das endgültige Modell auf neuen Daten?}
\]
\end{conceptbox}

\section{Typische Datenaufteilung}

Eine einfache Aufteilung ist zum Beispiel:

\[
    60\% \text{ Training},
    \qquad
    20\% \text{ Validierung},
    \qquad
    20\% \text{ Test}.
\]

Oder:

\[
    80\% \text{ Training},
    \qquad
    10\% \text{ Validierung},
    \qquad
    10\% \text{ Test}.
\]

Die genaue Aufteilung hängt von der Datenmenge ab.

\begin{notebox}
Bei sehr großen Datensätzen reicht oft ein kleiner Testanteil.
Bei kleinen Datensätzen möchte man möglichst viele Daten fürs Training nutzen.
Dann kann Kreuzvalidierung hilfreich sein.
\end{notebox}

\section{Kreuzvalidierung}

Bei kleinen Datensätzen ist eine einzelne Validierungsaufteilung oft instabil.
Das Ergebnis kann stark davon abhängen, welche Datenpunkte zufällig in Training oder Validierung landen.
Eine Lösung ist Kreuzvalidierung (cross-validation).

\begin{definitionbox}
Bei \(K\)-facher Kreuzvalidierung (\(K\)-fold cross-validation) wird der Datensatz in \(K\) Teile, sogenannte Folds, aufgeteilt.

Das Modell wird \(K\)-mal trainiert.
Jedes Mal wird ein anderer Fold zur Validierung verwendet und die übrigen \(K-1\) Folds zum Training.
\end{definitionbox}

Der Kreuzvalidierungsfehler ist der Durchschnitt der Validierungsfehler:

\[
    E_{\mathrm{CV}}
    =
    \frac{1}{K}
    \sum_{k=1}^{K}
    E_{\mathrm{val}}^{(k)}.
\]

Dabei ist \(E_{\mathrm{val}}^{(k)}\) der Validierungsfehler im \(k\)-ten Durchlauf.

\begin{conceptbox}
Kreuzvalidierung nutzt die Daten effizienter als eine einzelne Validierungsaufteilung.

Jeder Datenpunkt wird einmal zur Validierung und \(K-1\)-mal zum Training verwendet.
\end{conceptbox}

\section{Leave-One-Out-Kreuzvalidierung}

Ein Spezialfall ist Leave-One-Out-Kreuzvalidierung (leave-one-out cross-validation).

Dabei ist

\[
    K=N.
\]

Das bedeutet:
Jeder Fold enthält genau einen Datenpunkt.

Für jeden Durchlauf wird auf \(N-1\) Datenpunkten trainiert und auf dem übrig gebliebenen Datenpunkt validiert.

\begin{definitionbox}
Leave-One-Out-Kreuzvalidierung ist \(K\)-fold cross-validation mit

\[
    K=N.
\]
\end{definitionbox}

\begin{notebox}
Leave-One-Out nutzt die Daten sehr effizient, kann aber rechnerisch teuer sein, weil das Modell \(N\)-mal trainiert werden muss.
\end{notebox}

\section{Stratifizierte Aufteilung}

Bei Klassifikationsproblemen mit unausgeglichenen Klassen ist es wichtig, dass Training, Validierung und Test ähnliche Klassenanteile haben.

\begin{definitionbox}
Stratifizierte Aufteilung (stratified split) bedeutet, dass die Klassenverteilung in Training, Validierung und Test ungefähr gleich bleibt.
\end{definitionbox}

Beispiel:
Wenn im gesamten Datensatz

\[
    90\%
\]

der Datenpunkte Klasse \(0\) und

\[
    10\%
\]

Klasse \(1\) sind, sollten Training, Validierung und Test ungefähr dieselben Anteile haben.

\begin{examtipbox}
Bei unausgeglichenen Klassen sollte man nicht rein zufällig ohne Kontrolle splitten.
Sonst kann es passieren, dass eine seltene Klasse im Validierungs- oder Testdatensatz kaum vertreten ist.
\end{examtipbox}

\section{Data Leakage}

Data Leakage ist ein häufiger und gefährlicher Fehler.
Es bedeutet, dass Information aus Validierungs- oder Testdaten indirekt ins Training gelangt.

\begin{definitionbox}
Data Leakage bedeutet, dass ein Modell während des Trainings oder der Modellselektion Information verwendet, die eigentlich erst in Validierungs- oder Testdaten verfügbar sein sollte.
\end{definitionbox}

Beispiele:

\begin{itemize}
    \item Normalisierung mit Mittelwert und Standardabweichung des gesamten Datensatzes statt nur der Trainingsdaten.
    \item Feature-Auswahl auf allen Daten vor dem Split.
    \item Hyperparameterwahl anhand des Testfehlers.
    \item Duplikate desselben Falls in Training und Test.
    \item Zeitreihendaten zufällig mischen, obwohl zukünftige Daten nicht zur Vergangenheit verfügbar sein sollten.
\end{itemize}

\begin{conceptbox}
Data Leakage führt zu zu optimistischen Ergebnissen.

Das Modell wirkt besser, als es auf wirklich neuen Daten wäre.
\end{conceptbox}

\section{Vorverarbeitung und Splits}

Viele Vorverarbeitungsschritte müssen nur auf den Trainingsdaten gefittet werden.

Beispiel Standardisierung:

\[
    x_j'
    =
    \frac{x_j-\mu_j}{\sigma_j}.
\]

Dabei müssen \(\mu_j\) und \(\sigma_j\) nur aus den Trainingsdaten berechnet werden:

\[
    \mu_j
    =
    \frac{1}{N_{\mathrm{train}}}
    \sum_{n\in\mathcal{D}_{\mathrm{train}}}
    x_{nj},
\]

\[
    \sigma_j^2
    =
    \frac{1}{N_{\mathrm{train}}}
    \sum_{n\in\mathcal{D}_{\mathrm{train}}}
    (x_{nj}-\mu_j)^2.
\]

Dann werden dieselben Werte auf Validierungs- und Testdaten angewendet.

\begin{examtipbox}
Regel:

\[
    \text{Fit preprocessing on training data only.}
\]

Dann:

\[
    \text{Apply same transformation to validation and test data.}
\]
\end{examtipbox}

\section{Bias-Varianz-Trade-off}

Der Bias-Varianz-Trade-off wurde bereits bei Schätzern eingeführt.
Bei Generalisierung ist er besonders wichtig.

Informell:

\[
    \text{Bias}
    =
    \text{Fehler durch zu einfache Modellannahmen}.
\]

\[
    \text{Varianz}
    =
    \text{Empfindlichkeit gegenüber Trainingsdaten}.
\]

\begin{conceptbox}
Hoher Bias führt häufig zu Unteranpassung.

Hohe Varianz führt häufig zu Überanpassung.
\end{conceptbox}

\subsection{Hoher Bias}

Ein Modell mit hohem Bias ist zu starr.
Es kann die wahre Struktur nicht gut approximieren.

Typisch:

\[
    E_{\mathrm{train}}
    \text{ hoch},
    \qquad
    E_{\mathrm{test}}
    \text{ hoch}.
\]

\subsection{Hohe Varianz}

Ein Modell mit hoher Varianz reagiert stark auf die konkrete Trainingsmenge.
Es passt sich zufälligen Details an.

Typisch:

\[
    E_{\mathrm{train}}
    \text{ niedrig},
    \qquad
    E_{\mathrm{test}}
    \text{ hoch}.
\]

\section{Bias-Varianz-Zerlegung für Regression}

Für quadratischen Fehler kann man den erwarteten Vorhersagefehler in Bias, Varianz und Rauschen zerlegen.

Sei die wahre Beziehung:

\[
    t
    =
    f(\vec{x})
    +
    \epsilon,
\]

mit

\[
    \mathbb{E}[\epsilon]=0
\]

und

\[
    \mathrm{Var}(\epsilon)=\sigma^2.
\]

Ein Lernalgorithmus trainiert auf einem zufälligen Datensatz \(\mathcal{D}\) ein Modell

\[
    y(\vec{x};\mathcal{D}).
\]

Der erwartete quadratische Fehler an einem festen Punkt \(\vec{x}\) kann geschrieben werden als:

\[
    \mathbb{E}_{\mathcal{D},t}
    \left[
        (y(\vec{x};\mathcal{D})-t)^2
    \right]
    =
    \underbrace{
        \left(
            \mathbb{E}_{\mathcal{D}}[y(\vec{x};\mathcal{D})]
            -
            f(\vec{x})
        \right)^2
    }_{\text{Bias}^2}
    +
    \underbrace{
        \mathbb{E}_{\mathcal{D}}
        \left[
            \left(
                y(\vec{x};\mathcal{D})
                -
                \mathbb{E}_{\mathcal{D}}[y(\vec{x};\mathcal{D})]
            \right)^2
        \right]
    }_{\text{Varianz}}
    +
    \underbrace{\sigma^2}_{\text{Rauschen}}.
\]

\begin{definitionbox}
Für quadratischen Fehler zerfällt der erwartete Fehler in:

\[
    \text{erwarteter Fehler}
    =
    \text{Bias}^2
    +
    \text{Varianz}
    +
    \text{Rauschen}.
\]
\end{definitionbox}

\begin{notebox}
Das Rauschen \(\sigma^2\) ist irreduzibel.
Auch das beste Modell kann zufälliges Messrauschen nicht vorhersagen.
\end{notebox}

\section{Regularisierung}

Regularisierung (regularization) ist eine zentrale Methode gegen Overfitting.
Man bestraft zu komplexe Modelle oder große Parameterwerte.

Ein typischer regularisierter Verlust ist:

\[
    E_{\mathrm{reg}}(\theta)
    =
    E_{\mathrm{data}}(\theta)
    +
    \lambda E_{\mathrm{regterm}}(\theta).
\]

Dabei steuert \(\lambda\) die Stärke der Regularisierung.

\begin{definitionbox}
Regularisierung (regularization) fügt dem Trainingsziel einen Zusatzterm hinzu, der unerwünschte Modellkomplexität bestraft.
\end{definitionbox}

\subsection{\(L_2\)-Regularisierung}

Bei \(L_2\)-Regularisierung verwendet man:

\[
    E_{\mathrm{regterm}}(\vec{w})
    =
    \frac{1}{2}
    \vec{w}^{\top}\vec{w}.
\]

Der Verlust ist:

\[
    E_{\mathrm{reg}}(\vec{w})
    =
    E_{\mathrm{data}}(\vec{w})
    +
    \frac{\lambda}{2}
    \vec{w}^{\top}\vec{w}.
\]

\begin{conceptbox}
\(L_2\)-Regularisierung bestraft große Gewichte.

Sie führt oft zu glatteren, stabileren Modellen.
\end{conceptbox}

\subsection{\(L_1\)-Regularisierung}

Bei \(L_1\)-Regularisierung verwendet man:

\[
    E_{\mathrm{regterm}}(\vec{w})
    =
    \sum_{j}
    |w_j|.
\]

Der Verlust ist:

\[
    E_{\mathrm{reg}}(\vec{w})
    =
    E_{\mathrm{data}}(\vec{w})
    +
    \lambda
    \sum_{j}
    |w_j|.
\]

\begin{conceptbox}
\(L_1\)-Regularisierung fördert sparse Lösungen.

Das bedeutet:
Viele Gewichte werden exakt oder ungefähr \(0\).
\end{conceptbox}

\section{Regularisierung als MAP}

Regularisierung kann probabilistisch als Maximum-a-posteriori-Schätzung (maximum a-posteriori estimation) interpretiert werden.

Bei MAP maximieren wir:

\[
    p(\theta\mid\mathcal{D})
    \propto
    p(\mathcal{D}\mid\theta)p(\theta).
\]

Äquivalent minimieren wir den negativen Log-Posterior:

\[
    -\log p(\theta\mid\mathcal{D})
    =
    -\log p(\mathcal{D}\mid\theta)
    -
    \log p(\theta)
    +
    \text{const}.
\]

Der Datenfehler ist:

\[
    E_{\mathrm{data}}(\theta)
    =
    -\log p(\mathcal{D}\mid\theta).
\]

Der Regularisierungsterm kommt aus dem Prior:

\[
    E_{\mathrm{regterm}}(\theta)
    =
    -\log p(\theta).
\]

\begin{conceptbox}
Regularisierung entspricht einem Prior auf Parameter.

\[
    \text{Regularisierung}
    \quad
    \Longleftrightarrow
    \quad
    \text{Prior}.
\]
\end{conceptbox}

Ein Gauß-Prior führt zu \(L_2\)-Regularisierung.
Ein Laplace-Prior führt zu \(L_1\)-Regularisierung.

\section{Early Stopping}

Early Stopping ist eine Regularisierungsmethode für iterative Trainingsverfahren.

Beim Training sinkt der Trainingsfehler meistens weiter.
Der Validierungsfehler sinkt zunächst ebenfalls, kann aber später wieder steigen.

\[
    E_{\mathrm{train}}
    \downarrow
\]

während

\[
    E_{\mathrm{val}}
\]

nach einer gewissen Zeit wieder steigen kann.

\begin{definitionbox}
Early Stopping bedeutet, das Training zu stoppen, wenn der Validierungsfehler nicht mehr besser wird.
\end{definitionbox}

\begin{conceptbox}
Early Stopping verhindert, dass das Modell zu lange an die Trainingsdaten angepasst wird.

Es wirkt daher als Regularisierung.
\end{conceptbox}

Typisches Vorgehen:

\begin{enumerate}
    \item Trainiere iterativ.
    \item Berechne nach jeder Epoche den Validierungsfehler.
    \item Speichere das Modell mit dem besten Validierungsfehler.
    \item Stoppe, wenn sich der Validierungsfehler über mehrere Epochen nicht verbessert.
\end{enumerate}

\section{Lernkurven}

Lernkurven (learning curves) zeigen Trainings- und Validierungsfehler als Funktion der Trainingsmenge oder der Trainingszeit.

\subsection{Fehler über Trainingszeit}

Man zeichnet:

\[
    E_{\mathrm{train}}(\tau)
\]

und

\[
    E_{\mathrm{val}}(\tau)
\]

gegen die Epoche \(\tau\).

Typisches Overfitting-Muster:

\[
    E_{\mathrm{train}}(\tau)
    \text{ sinkt weiter},
\]

aber

\[
    E_{\mathrm{val}}(\tau)
    \text{ steigt wieder}.
\]

\subsection{Fehler über Datenmenge}

Man trainiert mit verschiedenen Trainingsgrößen \(N\) und zeichnet:

\[
    E_{\mathrm{train}}(N)
\]

und

\[
    E_{\mathrm{val}}(N).
\]

\begin{conceptbox}
Lernkurven helfen zu diagnostizieren, ob ein Modell eher unteranpasst oder überanpasst.
\end{conceptbox}

\section{Diagnose: Underfitting oder Overfitting}

\subsection{Fall 1: Beide Fehler hoch}

Wenn

\[
    E_{\mathrm{train}}
    \text{ hoch}
\]

und

\[
    E_{\mathrm{val}}
    \text{ hoch}
\]

sind, liegt wahrscheinlich Unteranpassung vor.

Mögliche Maßnahmen:

\begin{itemize}
    \item komplexeres Modell wählen,
    \item mehr Features verwenden,
    \item schwächere Regularisierung verwenden,
    \item länger trainieren,
    \item bessere Optimierung verwenden.
\end{itemize}

\subsection{Fall 2: Trainingsfehler niedrig, Validierungsfehler hoch}

Wenn

\[
    E_{\mathrm{train}}
    \text{ niedrig}
\]

aber

\[
    E_{\mathrm{val}}
    \text{ hoch}
\]

ist, liegt wahrscheinlich Überanpassung vor.

Mögliche Maßnahmen:

\begin{itemize}
    \item einfacheres Modell wählen,
    \item stärkere Regularisierung verwenden,
    \item mehr Trainingsdaten sammeln,
    \item Datenaugmentation verwenden,
    \item Early Stopping verwenden.
\end{itemize}

\subsection{Fall 3: Beide Fehler niedrig}

Wenn beide Fehler niedrig sind, generalisiert das Modell wahrscheinlich gut.
Trotzdem sollte am Ende ein unabhängiger Testdatensatz verwendet werden.

\section{Datenaugmentation}

Datenaugmentation (data augmentation) erzeugt zusätzliche Trainingsbeispiele durch transformationsbasierte Varianten vorhandener Daten.

Bei Bildern zum Beispiel:

\begin{itemize}
    \item kleine Rotationen,
    \item Spiegelungen,
    \item leichte Verschiebungen,
    \item zufällige Ausschnitte,
    \item Farbänderungen.
\end{itemize}

\begin{definitionbox}
Datenaugmentation (data augmentation) erweitert die Trainingsdaten durch sinnvolle Transformationen, die das Label erhalten.
\end{definitionbox}

\begin{conceptbox}
Datenaugmentation kann Overfitting reduzieren, weil das Modell mehr Variation sieht.
\end{conceptbox}

Wichtig:
Die Transformationen müssen label-erhaltend sein.
Wenn eine Transformation das Label verändert, erzeugt man falsche Trainingsdaten.

\section{Ensembles}

Ein Ensemble kombiniert mehrere Modelle.
Die Idee ist, dass verschiedene Modelle unterschiedliche Fehler machen.
Durch Mittelung oder Mehrheitsentscheidung kann die Gesamtleistung besser werden.

Für Regression kann man Vorhersagen mitteln:

\[
    y_{\mathrm{ens}}(\vec{x})
    =
    \frac{1}{M}
    \sum_{m=1}^{M}
    y_m(\vec{x}).
\]

Für Klassifikation kann man Wahrscheinlichkeiten mitteln:

\[
    p_{\mathrm{ens}}(c\mid\vec{x})
    =
    \frac{1}{M}
    \sum_{m=1}^{M}
    p_m(c\mid\vec{x}).
\]

\begin{definitionbox}
Ein Ensemble ist eine Kombination mehrerer Modelle zur Verbesserung von Stabilität und Generalisierung.
\end{definitionbox}

\begin{notebox}
Ensembles reduzieren oft Varianz, sind aber rechnerisch teurer und schwerer zu interpretieren.
\end{notebox}

\section{Modellselektion}

Modellselektion bedeutet, aus mehreren Kandidaten ein Modell auszuwählen.

Beispiele:

\begin{itemize}
    \item Polynomgrad \(M\) wählen,
    \item Regularisierung \(\lambda\) wählen,
    \item Lernrate \(\eta\) wählen,
    \item Netzwerkarchitektur wählen,
    \item Aktivierungsfunktion wählen.
\end{itemize}

Für jeden Kandidaten trainiert man auf den Trainingsdaten und bewertet auf den Validierungsdaten.

\[
    \lambda^*
    =
    \arg\min_{\lambda}
    E_{\mathrm{val}}(\lambda).
\]

\begin{definitionbox}
Hyperparameteroptimierung (hyperparameter optimization) ist die Suche nach Hyperparametern, die den Validierungsfehler minimieren.
\end{definitionbox}

\section{Grid Search und Random Search}

\subsection{Grid Search}

Bei Grid Search definiert man eine Menge möglicher Werte.
Zum Beispiel:

\[
    \lambda \in \{10^{-4},10^{-3},10^{-2},10^{-1},1\}.
\]

Dann trainiert man für jeden Wert ein Modell und wählt den besten Validierungsfehler.

\subsection{Random Search}

Bei Random Search werden Hyperparameter zufällig aus vorgegebenen Bereichen gezogen.
Das kann effizienter sein, wenn nur wenige Hyperparameter wirklich wichtig sind.

\begin{notebox}
Für Hyperparameter wie \(\lambda\) sucht man oft auf logarithmischer Skala, weil Unterschiede über Größenordnungen wichtig sind.
\end{notebox}

\section{Nested Cross-Validation}

Wenn sowohl Modellselektion als auch Leistungsbewertung mit Kreuzvalidierung gemacht werden sollen, verwendet man verschachtelte Kreuzvalidierung (nested cross-validation).

Es gibt:

\begin{itemize}
    \item eine äußere Kreuzvalidierung zur Bewertung,
    \item eine innere Kreuzvalidierung zur Hyperparameterwahl.
\end{itemize}

\begin{definitionbox}
Nested Cross-Validation trennt Modellselektion und Leistungsbewertung sauber durch zwei verschachtelte Validierungsschleifen.
\end{definitionbox}

\begin{conceptbox}
Nested Cross-Validation verhindert, dass die Leistungsbewertung durch Hyperparameterwahl zu optimistisch wird.
\end{conceptbox}

\section{Verteilungsshift}

Alle bisherigen Aussagen setzen voraus, dass Trainings- und Testdaten aus derselben Verteilung stammen:

\[
    p_{\mathrm{train}}(\vec{x},t)
    =
    p_{\mathrm{test}}(\vec{x},t).
\]

In der Praxis kann das verletzt sein.

\begin{definitionbox}
Verteilungsshift (distribution shift) bedeutet, dass sich die Datenverteilung zwischen Training und Anwendung unterscheidet.
\end{definitionbox}

Beispiele:

\begin{itemize}
    \item Ein medizinisches Modell wird in einem anderen Krankenhaus verwendet.
    \item Ein Spamfilter wird nach mehreren Jahren verwendet, obwohl sich Spam verändert hat.
    \item Ein Bildmodell wird auf Bildern mit anderer Kameraqualität getestet.
    \item Trainingsdaten stammen aus Simulation, Testdaten aus der realen Welt.
\end{itemize}

\begin{notebox}
Ein niedriger Testfehler ist nur aussagekräftig, wenn die Testdaten die spätere Anwendung gut repräsentieren.
\end{notebox}

\section{Metriken für Regression}

Bei Regression sind typische Metriken:

\subsection{Mean Squared Error}

\[
    \mathrm{MSE}
    =
    \frac{1}{N}
    \sum_{n=1}^{N}
    (y_n-t_n)^2.
\]

\subsection{Root Mean Squared Error}

\[
    \mathrm{RMSE}
    =
    \sqrt{
        \frac{1}{N}
        \sum_{n=1}^{N}
        (y_n-t_n)^2
    }.
\]

\subsection{Mean Absolute Error}

\[
    \mathrm{MAE}
    =
    \frac{1}{N}
    \sum_{n=1}^{N}
    |y_n-t_n|.
\]

\begin{conceptbox}
MSE bestraft große Fehler stärker als MAE, weil Fehler quadriert werden.

MAE ist robuster gegenüber Ausreißern.
\end{conceptbox}

\section{Metriken für Klassifikation}

Bei Klassifikation sind typische Metriken:

\subsection{Accuracy}

\[
    \mathrm{accuracy}
    =
    \frac{1}{N}
    \sum_{n=1}^{N}
    \mathbb{1}(\hat{t}_n=t_n).
\]

\subsection{Precision}

\[
    \mathrm{precision}
    =
    \frac{\mathrm{TP}}{\mathrm{TP}+\mathrm{FP}}.
\]

\subsection{Recall}

\[
    \mathrm{recall}
    =
    \frac{\mathrm{TP}}{\mathrm{TP}+\mathrm{FN}}.
\]

\subsection{F1-Score}

\[
    F_1
    =
    2
    \frac{
        \mathrm{precision}\cdot\mathrm{recall}
    }{
        \mathrm{precision}+\mathrm{recall}
    }.
\]

\begin{examtipbox}
Bei unausgeglichenen Klassen ist Accuracy oft nicht ausreichend.
Dann sind Precision, Recall, F1-Score oder ROC-AUC oft aussagekräftiger.
\end{examtipbox}

\section{Unsicherheit der Testschätzung}

Der Testfehler ist selbst nur ein Schätzer des wahren Generalisierungsfehlers.
Wenn der Testdatensatz klein ist, kann diese Schätzung stark schwanken.

Für Accuracy kann man grob modellieren:

\[
    \hat{p}
    =
    \frac{k}{N},
\]

wobei \(k\) die Anzahl korrekter Vorhersagen ist.
Die Standardabweichung dieses Anteils ist ungefähr:

\[
    \sqrt{
        \frac{\hat{p}(1-\hat{p})}{N}
    }.
\]

\begin{conceptbox}
Ein Testfehler auf wenigen Testpunkten ist unsicher.

Mehr Testdaten liefern eine stabilere Schätzung der echten Generalisierung.
\end{conceptbox}

\section{Häufige Fehler}

\begin{examtipbox}
Typische Fehler bei Generalisierung und Validierung:

\begin{enumerate}
    \item Trainingsfehler mit Generalisierungsfehler verwechseln.
    \item Testdaten zur Hyperparameterwahl verwenden.
    \item Vorverarbeitung auf dem gesamten Datensatz fitten.
    \item Overfitting nur an großem Trainingsfehler erkennen wollen.
    \item Underfitting und Overfitting verwechseln.
    \item Bei unausgeglichenen Klassen nur Accuracy betrachten.
    \item Kreuzvalidierung falsch durchführen, sodass Datenleckage entsteht.
    \item Bei Zeitreihen zufällig splitten, obwohl die zeitliche Reihenfolge wichtig ist.
    \item Den Validierungsfehler mehrfach optimieren und danach so tun, als wäre er ein unabhängiger Testfehler.
    \item Ein Modell auf einem Testdatensatz bewerten, der die spätere Anwendung nicht repräsentiert.
\end{enumerate}
\end{examtipbox}

\section{Mini-Aufgaben}

\subsection{Aufgabe 1: Trainings- und Testfehler interpretieren}

Ein Modell hat:

\[
    E_{\mathrm{train}} = 0{,}01
\]

und

\[
    E_{\mathrm{test}} = 0{,}50.
\]

Was ist wahrscheinlich passiert?

\begin{examplebox}
Der Trainingsfehler ist sehr klein, aber der Testfehler ist groß.
Das spricht für Überanpassung (overfitting).

Das Modell hat die Trainingsdaten sehr gut gelernt, generalisiert aber schlecht auf neue Daten.
\end{examplebox}

\subsection{Aufgabe 2: Underfitting erkennen}

Ein Modell hat:

\[
    E_{\mathrm{train}} = 0{,}45
\]

und

\[
    E_{\mathrm{val}} = 0{,}48.
\]

Was ist wahrscheinlich passiert?

\begin{examplebox}
Beide Fehler sind hoch.
Das spricht für Unteranpassung (underfitting).

Das Modell ist vermutlich zu einfach oder wurde nicht ausreichend trainiert.
\end{examplebox}

\subsection{Aufgabe 3: Kreuzvalidierung}

Ein \(5\)-fold cross-validation Verfahren liefert Validierungsfehler:

\[
    0{,}20,
    \quad
    0{,}25,
    \quad
    0{,}22,
    \quad
    0{,}23,
    \quad
    0{,}30.
\]

Berechne den Kreuzvalidierungsfehler.

\begin{examplebox}
Der Kreuzvalidierungsfehler ist der Durchschnitt:

\[
    E_{\mathrm{CV}}
    =
    \frac{1}{5}
    (0{,}20+0{,}25+0{,}22+0{,}23+0{,}30).
\]

Die Summe ist:

\[
    0{,}20+0{,}25+0{,}22+0{,}23+0{,}30
    =
    1{,}20.
\]

Also:

\[
    E_{\mathrm{CV}}
    =
    \frac{1{,}20}{5}
    =
    0{,}24.
\]
\end{examplebox}

\subsection{Aufgabe 4: Accuracy}

Ein Klassifikator macht auf \(100\) Testpunkten \(87\) korrekte Vorhersagen.
Berechne die Accuracy.

\begin{examplebox}
Die Accuracy ist:

\[
    \mathrm{accuracy}
    =
    \frac{87}{100}
    =
    0{,}87.
\]

Also:

\[
    \mathrm{accuracy}
    =
    87\%.
\]
\end{examplebox}

\subsection{Aufgabe 5: Precision und Recall}

Gegeben seien:

\[
    \mathrm{TP}=30,
    \qquad
    \mathrm{FP}=10,
    \qquad
    \mathrm{FN}=20.
\]

Berechne Precision und Recall.

\begin{examplebox}
Precision:

\[
    \mathrm{precision}
    =
    \frac{\mathrm{TP}}{\mathrm{TP}+\mathrm{FP}}
    =
    \frac{30}{30+10}
    =
    \frac{30}{40}
    =
    0{,}75.
\]

Recall:

\[
    \mathrm{recall}
    =
    \frac{\mathrm{TP}}{\mathrm{TP}+\mathrm{FN}}
    =
    \frac{30}{30+20}
    =
    \frac{30}{50}
    =
    0{,}60.
\]
\end{examplebox}

\section{Zusammenfassung}

\begin{itemize}
    \item Generalisierung (generalization) bedeutet gute Leistung auf neuen Daten.
    \item Der Trainingsfehler misst Leistung auf Trainingsdaten.
    \item Der Validierungsfehler dient zur Modell- und Hyperparameterwahl.
    \item Der Testfehler dient zur finalen Leistungsabschätzung.
    \item Überanpassung (overfitting) bedeutet niedriger Trainingsfehler, aber hoher Fehler auf neuen Daten.
    \item Unteranpassung (underfitting) bedeutet, dass das Modell zu einfach ist und selbst die Trainingsdaten schlecht beschreibt.
    \item Modellkomplexität beeinflusst den Bias-Varianz-Trade-off.
    \item Hyperparameter werden nicht direkt durch Training gelernt.
    \item Kreuzvalidierung nutzt Daten effizienter für Modellselektion.
    \item Testdaten dürfen nicht zur Hyperparameterwahl verwendet werden.
    \item Data Leakage führt zu zu optimistischen Ergebnissen.
    \item Regularisierung reduziert Overfitting durch Bestrafung zu komplexer Modelle.
    \item Early Stopping ist eine Regularisierungsmethode für iterative Trainingsverfahren.
    \item Lernkurven helfen, Underfitting und Overfitting zu diagnostizieren.
    \item Die passende Metrik hängt vom Problem ab.
\end{itemize}

\section{Typische Prüfungsfragen}

\begin{examtipbox}
Mögliche Prüfungsfragen zu diesem Kapitel:

\begin{enumerate}
    \item Was bedeutet Generalisierung?
    \item Warum reicht ein kleiner Trainingsfehler nicht aus?
    \item Was ist der Unterschied zwischen Trainings-, Validierungs- und Testdaten?
    \item Warum darf man Testdaten nicht zur Hyperparameterwahl verwenden?
    \item Was ist Overfitting?
    \item Was ist Underfitting?
    \item Wie erkennt man Overfitting an Trainings- und Validierungsfehler?
    \item Was ist Modellkomplexität?
    \item Was ist ein Hyperparameter?
    \item Was ist Modellselektion?
    \item Erkläre \(K\)-fold cross-validation.
    \item Was ist Leave-One-Out-Kreuzvalidierung?
    \item Was ist Data Leakage?
    \item Warum muss Standardisierung nur auf Trainingsdaten gefittet werden?
    \item Erkläre den Bias-Varianz-Trade-off.
    \item Was ist der irreduzible Fehler?
    \item Wie hilft Regularisierung gegen Overfitting?
    \item Wie hängt Regularisierung mit MAP zusammen?
    \item Was ist Early Stopping?
    \item Was zeigen Lernkurven?
    \item Wann ist Accuracy eine schlechte Metrik?
    \item Was ist der Unterschied zwischen Precision und Recall?
    \item Was ist Verteilungsshift?
    \item Warum ist ein Testfehler auf wenigen Daten unsicher?
\end{enumerate}
\end{examtipbox}